% ============ Solutions and hints (Part 10: §28, §29, §30) ============
% 原书 pp.296-297; 原书对 §27 的 Ex.27.1, 27.2、§31-§33 的习题未给解答, 编号按原书跳号.

\section*{\S\ref{sec:28}}

\noindent\textbf{28.1.}\quad Let $N$ be a $B$-module satisfying $m^\nu N = 0$, and $D\colon B \to N$ a
derivation over $A$; then $D$ induces a derivation $\bar{D}\colon B_0 = B/mB \to N/mN$.
If $B_0$ is 0-unramified over $k$ then $\bar{D} = 0$, so that $D(B) \subset mN$.
Proceeding as before gives $D(B) \subset m^2N, \ldots$ so that $D = 0$. The statement
about \'etale is just putting together those for smooth and unramified.

\noindent\textbf{28.2.}\quad Since some $D \in \Der(k)$ with $Da \ne 0$ can be extended to a derivation
of $A$, $a \notin A^p$. If $k'$ were a coefficient field containing $k$ then we would
have to have $a \in k'^p \subset A^p$. Also, $A$ is 0-smooth over $k$ because
$k[X]$ is.

\section*{\S\ref{sec:29}}

\noindent\textbf{29.1.}\quad Suppose that $C = R[[y]]$ with $R$ a DVR; if we let $u$ be a
uniformising element of $R$ then $pR$ is a power of $uR$, so that $pC$ has the
single prime divisor $uC$. However, in fact, in our case
$C/pC = (B/pB)[X]/(X(X + y))$, so that $pC$ has the two prime divisors $(p, x)$
and $(p, x + y)$.

\noindent\textbf{29.2.}\quad $R$ is $pR$-\'etale over $\mathbb{Z}_{(p)}$ (see Ex.~\hyperref[ex:28.1-ex]{28.1}).

\section*{\S\ref{sec:30}}

\noindent\textbf{30.1.}\quad $\varphi$ is the composite of $E_t\colon A \to A[[t]]$ and of the map
$A[[t]] \to A$ obtained by substituting $-x$ for $t$. Since $\varphi(x) = 0$ we have
$xA \subset \Ker \varphi$. For any $a \in A$ we can write $\varphi(a) = a + xb$, so that
$\varphi(\varphi(a)) = \varphi(a)$, and therefore $C \cap xA = (0)$ and $A = C + xA$,
therefore $A = C[x]$. Now $E_t(x) = x + t$, and it is easy to see that this is a
non-zero-divisor of $A[[t]]$, so that $x$ is a non-zero-divisor of $A$. If
$c_\nu x^\nu + c_{\nu+1}x^{\nu+1} + \cdots = 0$ with $c_i \in C$ then dividing by $x^\nu$ we get
$c_\nu \in C \cap xA = (0)$, so that $x$ is analytically independent over $C$.

\noindent\textbf{30.2.}\quad No. If $A$ is 0-smooth over $k'$ then so is the field of fractions $L$
of $A$, so $L/k'$ is separable, hence also $k/k'$, and this is not the case.

\noindent\textbf{30.3.}\quad (1) $\Rightarrow$ (3) is easy using Theorem~\hyperref[thm:28.7]{28.7}. If
$[k : k^p] = \infty$ then there are counter-examples to (3) $\Rightarrow$ (1) (see [C1],
(22.7.7)).
